\appendix

%!TEX root = ../main.tex
\begin{table}[t]
  \centering
  \caption{Operator Space}
  \vspace{-1em}
  \resizebox{0.99\linewidth}{!}{
  \begin{tabular}{|c|l|l|} 
    \hline
    \textbf{Type}                & \textbf{Operator}      & \textbf{Typical Parameters}                   \\
    \hline \hline   
    \multirow{7}{*}{\makecell{Data \\ Cleaning}}    
                                          & StandardizeString      & column, func                            \\
                                          & ErrorDetection         & column, func                            \\
										  & DropNA                 & subset, how                                   \\
                                          & MVI & column, mode                            \\
                                          & OutlierDetection       & column, mode, action                    \\
                                          & StandardizeDatetime    & column, date\_format                    \\
                                          & Deduplicate            & subset, keep                                  \\ \hline \hline 
    \multirow{7}{*}{\makecell{Column \\ Transformation}} & AddNewColumn           & func                                          \\
                                          & Cast                   & column, dtype                                 \\
                                          & Explode                & column, split\_comma                          \\
                                          & RenameColumn           & rename\_map                                   \\
                                          & DropColumn             & drop\_columns                                 \\
                                          & SplitColumn            & src\_column, func                            \\
                                          & Concatenate            & columns, func                    \\ \hline \hline 
    \multirow{15}{*}{\makecell{Table \\ Transformation}}    & Join                   & left, right, on, how                          \\
                                          & Union                  & tables, how                                   \\
                                          & Append                 & src\_table, appended\_table  			\\ 
                                          & Pivot                  & index, columns, values               \\
                                          & Transpose              & -                                             \\
                                          & Stack                  & id\_vars, value\_vars \\
                                          & WideToLong             & subnames, i, j, sep, suffix                   \\
                                          & Subtitle               & title                                         \\
                                          & Sort                   & by, ascending                                 \\
                                          & GroupBy                & by, agg                                       \\
                                          & TopK                   & k                                             \\ 
											& Filter                 & func                            				\\
                                          & SelectColumn           & columns                                       \\
                                          & Count                  & table\_name                                   \\
                                          & CalculateStatistic     & tables, var\_name, func                 \\ \hline \hline 
    \multirow{2}{*}{Other}                & CodeGen         & tables, code                            \\
                                          & Terminate              & output\_tables \\ \hline                              
  \end{tabular}
  }
  \label{tbl:operator_space}
  % \vspace{-1em}
\end{table}

\section{Complete Operator Definitions}
\label{sec:appendix_operators}

In this section, we provide a comprehensive definition of the operator space $\mathcal{O}$ utilized in \sys. As outlined in Table~\ref{tbl:operator_space}, the operators are categorized into four groups: Data Cleaning, Column Transformation, Table Reconstruction, and Other. We detail the functionality and parameters for each operator as below.

\stitle{Operators for Data Cleaning.}
These operators focus on improving data quality by handling errors, missing values, and inconsistencies.

(1) $\mathtt{StandardizeString}$.
    This function standardizes the string format within a specific column. It accepts two arguments: {\small\texttt{column}}, representing the target column name, and {\small\texttt{func}}, a lambda function defining the transformation logic. For example, to remove whitespace and convert the \term{City} column to title case, we set {\small\texttt{column}} to \term{City} and {\small\texttt{func}} to {\small\texttt{lambda x: x.strip().title()}}.

(2) $\mathtt{ErrorDetection}$.
    This function identifies erroneous values based on specific logic. It takes {\small\texttt{column}} and {\small\texttt{func}} as arguments, where {\small\texttt{func}} returns a boolean indicating validity. For instance, to detect invalid email addresses in an \term{Email} column, {\small\texttt{func}} could be a regex match function {\small\texttt{lambda x: not re.match(r"[\^{}@]+@[\^{}@]+\.[\^{}@]+", x)}}.

(3) $\mathtt{DropNA}$.
    This function removes rows containing missing values. It requires {\small\texttt{subset}}, a list of columns to inspect for null values, and {\small\texttt{how}}, which specifies the condition (\eg {\small\texttt{'any'}} or {\small\texttt{'all'}}). For example, to drop rows where both \term{Revenue} and \term{Cost} are missing, we set {\small\texttt{subset}} to {\small\texttt{['Revenue', 'Cost']}} and {\small\texttt{how}} to {\small\texttt{'all'}}.

(4) $\mathtt{MVI}$ (Missing Value Imputation).
    This function fills missing values using a specified strategy. It takes {\small\texttt{column}} and {\small\texttt{mode}} as arguments. The {\small\texttt{mode}} can be a statistical method (\eg {\small\texttt{'mean'}}, {\small\texttt{'median'}}) or a constant value. For example, to fill missing values in the \term{Age} column with the average age, we set {\small\texttt{column}} to \term{Age} and {\small\texttt{mode}} to {\small\texttt{'mean'}}.

(5) $\mathtt{OutlierDetection}$.
    This function handles statistical outliers. It accepts {\small\texttt{column}}, {\small\texttt{mode}} (\eg {\small\texttt{'z-score'}}, {\small\texttt{'iqr'}}), and {\small\texttt{action}} (\eg {\small\texttt{'clip'}}, {\small\texttt{'drop'}}). For instance, to cap values in the \term{Price} column that exceed 3 standard deviations, we set {\small\texttt{mode}} to {\small\texttt{'z-score'}} and {\small\texttt{action}} to {\small\texttt{'clip'}}.

(6) $\mathtt{StandardizeDatetime}$.
    This function converts a column to a standard datetime object. It requires {\small\texttt{column}} and {\small\texttt{date\_format}}. For example, to parse dates like "2023/12/25" in the \term{Date} column, we set {\small\texttt{date\_format}} to {\small\texttt{"\%Y/\%m/\%d"}}.

(7) $\mathtt{Deduplicate}$.
    This function removes duplicate rows. It takes {\small\texttt{subset}} (columns to consider for identifying duplicates) and {\small\texttt{keep}} (which duplicate to retain, \eg {\small\texttt{'first'}}, {\small\texttt{'last'}}, or {\small\texttt{'none'}}).

\stitle{Operators for Column Transformation.}
These operators modify the schema by adding, removing, or restructuring columns.

(1) $\mathtt{AddNewColumn}$.
    This function creates a new column based on row-wise calculation. It takes a single argument {\small\texttt{func}}, a lambda function accepting a row dictionary. For example, to create a \term{Profit} column, we use {\small\texttt{lambda row: row['Revenue'] - row['Cost']}}.

(2) $\mathtt{Cast}$.
    This function converts the data type of a column. It arguments are {\small\texttt{column}} and {\small\texttt{dtype}} (\eg {\small\texttt{'int'}}, {\small\texttt{'float'}}, {\small\texttt{'str'}}). For example, to convert a \term{Year} column stored as text to integers, we set {\small\texttt{dtype}} to {\small\texttt{'int'}}.

(3) $\mathtt{Explode}$.
    This function transforms each element of a list-like string into a row, replicating index values. It takes {\small\texttt{column}} and {\small\texttt{split\_comma}} (the separator). As shown in Example~\ref{exam:agentic_tree_based_reasoning}, to split "Sci-Fi/Action" in the \term{Genre} column into separate rows, we set {\small\texttt{split\_comma}} to {\small\texttt{'/'}}.

(4) $\mathtt{RenameColumn}$.
    This function renames existing columns. It accepts {\small\texttt{rename\_map}}, a dictionary mapping old names to new names. \eg {\small\texttt{\{'Yr': 'Year', 'Pop': 'Population'\}}}.

(5) $\mathtt{DropColumn}$.
    This function removes columns from the table. It takes {\small\texttt{drop\_columns}}, a list of column names to delete.

(6) $\mathtt{SplitColumn}$.
    This function splits a single column into multiple columns based on a pattern. It requires {\small\texttt{src\_column}} and {\small\texttt{func}}. For example, to extract the numeric rating from a string "8.5/10", we set {\small\texttt{func}} to {\small\texttt{lambda x: x.split('/')[0]}}.

(7) $\mathtt{Concatenate}$.
    This function combines multiple columns into one string column. It takes {\small\texttt{columns}} (list of source columns) and {\small\texttt{func}} (formatting logic). For example, to combine \term{First} and \term{Last} names, {\small\texttt{func}} could be {\small\texttt{lambda x: f"\{x['First']\} \{x['Last']\}"}}.

\stitle{Operators for Table Reconstruction.}
These operators restructure the table shape or aggregate data.

(1) $\mathtt{Join}$.
    This function merges two tables. Arguments include {\small\texttt{left}} and {\small\texttt{right}} (table names), {\small\texttt{on}} (join keys), and {\small\texttt{how}} (type of join: {\small\texttt{'inner'}}, {\small\texttt{'left'}}, {\small\texttt{'outer'}}).

(2) $\mathtt{Union}$ / $\mathtt{Append}$.
    These functions combine tables vertically. $\mathtt{Union}$ takes a list of {\small\texttt{tables}} and a {\small\texttt{how}} parameter (\eg set logic), while $\mathtt{Append}$ specifically adds {\small\texttt{appended\_table}} to the end of {\small\texttt{src\_table}}.

(3) $\mathtt{Pivot}$.
    This function reshapes data from long to wide format. It requires {\small\texttt{index}}, {\small\texttt{columns}} (keys to become new headers), and {\small\texttt{values}}. For example, pivoting a sales table to show months as columns.

(4) $\mathtt{Transpose}$.
    This function swaps rows and columns of the table. It takes no parameters.

(5) $\mathtt{Stack}$ / $\mathtt{WideToLong}$.
    These functions reshape data from wide to long format. $\mathtt{Stack}$ takes {\small\texttt{id\_vars}} (columns to keep) and {\small\texttt{value\_vars}} (columns to unpivot). $\mathtt{WideToLong}$ handles more complex patterns with suffix extraction defined by {\small\texttt{subnames}}, {\small\texttt{i}}, {\small\texttt{j}}, and {\small\texttt{sep}}.

(6) $\mathtt{Subtitle}$.
    This function is used to handle hierarchical headers often found in spreadsheets. It takes {\small\texttt{title}} to identify the main header row and propagates it to sub-headers.

(7) $\mathtt{Sort}$.
    This function sorts the table rows. It takes {\small\texttt{by}} (column names) and {\small\texttt{ascending}} (boolean).

(8) $\mathtt{GroupBy}$.
    This function groups data and applies aggregation. It requires {\small\texttt{by}} (grouping columns) and {\small\texttt{agg}} (a dictionary mapping columns to functions, \eg {\small\texttt{\{'Profit': 'mean'\}}}).

(9) $\mathtt{TopK}$.
    This function retains the top-$k$ rows based on a specific order. It takes the integer argument {\small\texttt{k}}.

(10) $\mathtt{Filter}$.
    This function filters rows based on a condition. It takes {\small\texttt{func}}, a lambda function returning a boolean for each row. \eg {\small\texttt{lambda row: row['Year'] > 2000}}.

(11) $\mathtt{SelectColumn}$.
    This function selects a subset of columns to keep, discarding the rest. It takes {\small\texttt{columns}} as a list.

(12) $\mathtt{Count}$.
    This function counts the number of records, optionally grouped by unique values. It takes {\small\texttt{table\_name}}.

(13) $\mathtt{CalculateStatistic}$.
    This function computes statistics across tables or complex aggregations. It takes {\small\texttt{tables}}, {\small\texttt{var\_name}}, and {\small\texttt{func}}.

\stitle{Other Operators.}
These operators handle special cases and flow control.

(1) $\mathtt{CodeGeneration}$.
    This function allows the agent to generate arbitrary Python code for tasks not covered by the predefined operators. It takes {\small\texttt{tables}} and the generated {\small\texttt{code}} string. This ensures \sys can handle long-tail requirements.

(2) $\mathtt{Terminate}$.
    This operator signals the end of the reasoning process. It takes {\small\texttt{output\_tables}} to specify the final result to be returned to the user.

In total, we define \textbf{31} distinct operators, providing comprehensive coverage for diverse data cleaning and transformation tasks.